\section{Архитектура на библиотеката}
\label{sec:architecture}

Архитектурата на библиотеката е организирана около ясно разграничение между логическия модел на данните, междинното представяне на заявките и backend-специфичния слой за изпълнение. Това разграничение е ключово за целта на разработката: промяна на базата данни да изисква минимална промяна в приложния код, а в идеалния случай --- единствено смяна на типа на конектора. За да бъде постигнато това, архитектурата е изградена от няколко слоя, всеки със собствена отговорност и със строго определени интерфейси.

В теоретичната глава беше показано, че унифицирането е възможно на равнището на логическия модел и query IR-а, но не и чрез игнориране на различията между backend-ите. Настоящата глава прави следващата стъпка: показва как тази теоретична постановка е превърната в конкретна архитектура. Тук библиотеката вече не се разглежда само като идея, а като организирана система от типове, contracts и execution path-ове.

\subsection{Архитектурни цели и ограничения}

Архитектурата е проектирана при едновременно действие на няколко ограничения. Първо, user-facing API трябва да остане в рамките на идиоматичен C++, без въвеждане на отделен DSL и без runtime reflection механизъм. Второ, query структурата трябва да бъде представима като \code{constexpr} обект, така че значима част от валидирането да се извършва от компилатора. Трето, backend слойът трябва да бъде разширяем без наследяване от общ виртуален интерфейс, тъй като подобен интерфейс би скрил различията между backend-ите и би изместил проверките към runtime.

Четвъртото ограничение е свързано със самия multi-backend характер на системата. Библиотеката трябва да може да поддържа SQL и non-SQL конектори без фалшиво уеднаквяване. Това означава, че архитектурата трябва едновременно да допуска общ API и да изразява ограниченията на конкретен backend по формален и проверим начин. Оттук естествено следва изборът на trait-базиран connector слой и capability механизъм.

\begin{figure}[H]
\centering
\begin{tikzpicture}[node distance=1.3cm and 1.4cm]
    \node[ormaccent] (app) {\textbf{Приложен слой}\\entity типове и query композиция};
    \node[ormblock, below=of app] (entity) {\textbf{Entity слой}\\\code{property}, \code{relationship}, \code{table\_name\_trait}};
    \node[ormblock, below=of entity] (ir) {\textbf{Query IR слой}\\\code{field}, \code{Rule}, \code{select\_query}, \code{insert\_query}};
    \node[ormblock, below=of ir] (connector) {\textbf{Connector слой}\\\code{connector\_trait<DB>} и capability tags};
    \node[ormblock, below=of connector] (driver) {\textbf{Driver / runtime слой}\\SQL/BSON/Redis/Cypher/CQL материализация};

    \draw[ormline] (app) -- (entity);
    \draw[ormline] (entity) -- (ir);
    \draw[ormline] (ir) -- (connector);
    \draw[ormline] (connector) -- (driver);
\end{tikzpicture}
\caption{Слоеста архитектура на библиотеката}
\label{fig:architecture_layers}
\end{figure}

Фигура~\ref{fig:architecture_layers} показва, че библиотеката следва последователна вертикална организация. Приложният код не комуникира директно с драйвера; той работи с C++ типове и query builder-и. Драйверът от своя страна не познава домейн семантиката, а само materialized формата, генерирана от connector слоя. Именно поради това архитектурата може да остане едновременно строга и надграждаема.

\subsection{Слоеста организация и разделяне на отговорностите}

Най-горният слой е приложният код, който дефинира entity типове, изгражда заявки чрез \code{select}, \code{insert}, \code{update} и \code{deleteq}, и ги подава на \code{db<DB>}. Следващият слой е логическият модел на библиотеката --- entity описанията и query IR. Третият слой е connector abstraction-ът, където \code{connector\_trait<DB>} материализира същата логика към backend-специфичен синтаксис и поведение. Най-долу стои конкретният драйвер, wire protocol или mock implementation.

Тази архитектура има важен практически ефект. Ако потребителят промени \code{SQLiteDB} с \code{MongoDB}, query IR остава логически същият, но конекторът го тълкува като filter/projection изпълнение вместо SQL рендиране. Това не означава, че всеки IR е допустим върху всеки backend; означава, че \emph{формата на абстракцията} е обща, а допустимостта се определя формално чрез capability механизма.

\begin{table}[H]
\centering
\small
\caption{Основни архитектурни слоеве и техните отговорности}
\label{tab:architecture_responsibilities}
\begin{tabularx}{\textwidth}{L{2.8cm}L{4.0cm}Y}
\toprule
\textbf{Слой} & \textbf{Основни артефакти} & \textbf{Отговорност} \\
\midrule
Приложен & entity типове, query builder повиквания & описва домейн логика и формулира операции в идиоматичен C++ \\
Логически & \code{property}, \code{relationship}, \code{field}, \code{Rule} & моделира данните и заявките независимо от конкретен backend \\
Connector & \code{connector\_trait<DB>}, capability tags & проверява допустимостта и материализира IR-а към backend поведение \\
Runtime/driver & конкретни connection handle-и и драйверни API-та & изпълнява материализираната заявка и връща резултати \\
\bottomrule
\end{tabularx}
\end{table}

Таблица~\ref{tab:architecture_responsibilities} показва, че библиотеката не използва ``голям обект'', който да прави всичко. Вместо това различните решения са локализирани в отделни слоеве. Това е важно както за надграждаемостта, така и за академичния анализ, защото позволява отделните твърдения да бъдат проследени до конкретни типове и модули.

\subsection{Entity слой като архитектурна основа}

Entity слоят е изграден върху \code{property<T, Name>}, \code{relationship<...>} и \code{table\_name\_trait<T>}. Той има двойна роля. От една страна, представя C++ домейн модела. От друга страна, служи като логическо описание на схемата или структурата на персистентните данни. Това означава, че архитектурата не прави отделна runtime registry от метаданни. Вместо това типовата система съдържа достатъчно информация, за да бъдат извлечени имена на полета, типове, връзки и целеви контейнер за съхранение.

Особено важен е фактът, че entity слоят е независим от вида база данни. Нито \code{property}, нито \code{relationship} са SQL-специфични. Те описват логически атрибути и отношения. Така се създава основата за сравнимо поведение на релационни и нерелационни backend-и. На архитектурно ниво именно това прави възможно библиотеката да твърди, че C++ моделът е първичният носител на логическата схема.

\subsection{Query IR като централен архитектурен компонент}

Query IR е централният архитектурен компонент. Той обединява field референции, предикати, placeholder-и и допълнителни клаузи като \code{order\_by}, \code{group\_by} и \code{limit}. Вместо да се работи с текст на заявка, библиотеката изгражда типизиран обект, чийто шаблонни параметри кодират логическата структура на операцията.

От архитектурна гледна точка това решение има няколко последствия. Първо, валидирането на структурата се извършва рано. Второ, query IR може да бъде анализиран от connector слоя без загуба на семантична информация. Трето, подготвените заявки могат да запазят не SQL низ, а самото междинно представяне, което е по-естествено за библиотека, ориентирана към C++ типовата система.

Четвърто, query IR-ът е единствената зона, в която се срещат entity описанието и по-късната backend материализация. В този смисъл той е архитектурният ``шарнир'' на цялата библиотека: достатъчно висок, за да остане storage-agnostic, и достатъчно конкретен, за да може да бъде изпълнен.

\subsection{Конекторен слой и формалният contract}

Формално connector слоят е реализиран чрез специализация на \code{connector\_trait<DB>}. Този подход е аналогичен на стандартни trait шаблони като \code{std::iterator\_traits}. Типът \code{DB} може да бъде прост wrapper над connection handle, а цялата ORM логика за него да остане в trait специализацията. По този начин библиотеката не изисква потребителят или авторът на конектора да наследява общ базов клас.

\code{connector\_trait<DB>} определя поне три ключови аспекта: \code{wire\_type}, \code{cursor\_type} и \code{execute(...)}. По желание специализацията декларира и capability tags като \code{supports\_joins}, \code{supports\_transactions}, \code{supports\_aggregation} или \code{supports\_concurrent\_execute}. Наличието или отсъствието на тези псевдо-маркери е достатъчно, за да може по-горният слой да наложи compile-time ограничения.

\begin{lstlisting}[language=C++,caption={Извадка от формалния connector contract},label={lst:connector_contract}]
template <typename DB>
concept is_connector = requires {
    typename connector_trait<DB>::template wire_type<int>::type;
    typename connector_trait<DB>::cursor_type;
};

template <typename DB, typename Cap>
inline constexpr bool has_capability =
    detail::capability_check<DB, Cap>::value;
\end{lstlisting}

Listing~\ref{lst:connector_contract} показва, че contract-ът е структурен, а не наследствен. Това е решаващ архитектурен избор. Ако съществуваше виртуален интерфейс с runtime dispatch, част от типово проверимата информация щеше да бъде изгубена. Вместо това библиотеката предпочита compile-time структура, при която изискванията са видими и за компилатора, и за автора на конектора.

\begin{figure}[H]
\centering
\begin{tikzpicture}[node distance=1.35cm and 1.25cm]
    \node[ormaccent] (query) {\textbf{Типизирана заявка}\\join, transaction, aggregation intent};
    \node[ormblock, below=of query] (db) {\code{db<DB>} фасада};
    \node[ormblock, below left=of db] (capok) {Capability налице\\заявката е допустима};
    \node[ormdanger, below right=of db] (capfail) {Capability липсва\\compile-time отказ};
    \node[ormblock, below=of capok] (exec) {\code{connector\_trait<DB>::execute}};

    \draw[ormline] (query) -- (db);
    \draw[ormline] (db) -- (capok);
    \draw[ormline] (db) -- (capfail);
    \draw[ormline] (capok) -- (exec);
\end{tikzpicture}
\caption{Capability gating между query IR-а и connector слоя}
\label{fig:capability_gating}
\end{figure}

Фигура~\ref{fig:capability_gating} е важна, защото показва къде архитектурата взема решение за допустимост. Това не е драйверът и не е runtime логът. Решението се взема в публичния API слой, който вече има достъп до query IR-а и до trait информацията за конектора.

\begin{table}[H]
\centering
\small
\caption{Роля на основните capability tags в архитектурата}
\label{tab:capability_tags}
\begin{tabularx}{\textwidth}{L{3.6cm}Y}
\toprule
\textbf{Capability} & \textbf{Архитектурен смисъл} \\
\midrule
\code{supports\_joins} & разрешава join-базираните query форми за съответния backend \\
\code{supports\_transactions} & позволява използване на \code{transaction\_guard} и \code{begin}/\code{commit}/\code{rollback} hooks \\
\code{supports\_aggregation} & обозначава поддръжка на aggregation операции в логическия query слой \\
\code{supports\_embedding} & прави експлицитна native поддръжка на вградени структури \\
\code{supports\_upsert} & маркира backend-и, при които upsert е част от естествения execution contract \\
\code{supports\_bulk\_insert} & обозначава наличието на по-ефективен път за масово записване \\
\bottomrule
\end{tabularx}
\end{table}

\subsection{Потребителският фасаден слой \code{db<DB>}}

Класът \code{db<DB>} е публичният фасаден слой. Той капсулира препратка към конкретния backend обект и предоставя уеднаквен API за изпълнение. На практика този слой има две роли. Първата е удобство за потребителя: библиотеката се използва чрез един и същ интерфейс независимо от backend-а. Втората е архитектурен филтър: точно тук се извършват compile-time проверки за capability и съвместимост.

Следователно \code{db<DB>} не е просто тънък wrapper над \code{execute}. Той е мястото, където библиотеката формализира границата между разрешено и неразрешено поведение за избрания backend. Точно този слой позволява съобщения за грешка от вида ``конекторът не поддържа \code{JOIN} операции'', вместо грешката да се появи по-късно като неясно runtime поведение.

\begin{lstlisting}[language=C++,caption={Извадка от фасадния слой \code{db<DB>}},label={lst:db_facade}]
template <typename DB>
    requires is_connector<DB>
class db
{
public:
    template <typename Query>
    auto operator<<(Query q)
    {
        if constexpr (is_select_query<Query>)
        {
            if constexpr (decltype(q.join_clauses())::size > 0)
            {
                static_assert(has_capability<DB, cap::supports_joins>,
                    "orm::db: this connector does not support JOIN operations.");
            }
        }
        return connector_trait<DB>::execute(*conn_, std::move(q));
    }
};
\end{lstlisting}

Listing~\ref{lst:db_facade} показва защо \code{db<DB>} е повече от удобна фасада. Той е формалното място, където query структурата и connector capability информацията се срещат. Именно тук compile-time архитектурата се превръща в реално ограничение върху API употребата.

\subsection{Миграции като архитектурно разширение}

Архитектурата включва и прототипен migration слой чрез \code{migrate<DB>}, \code{live\_schema}, \code{entity\_meta} и DDL операции като \code{create\_table\_op}, \code{add\_column\_op}, \code{drop\_column\_op} и \code{alter\_column\_type\_op}. Този слой е особено важен от гледна точка на тезата, защото прави явна връзката между C++ модела и реалната схема на съхранение.

Миграционният слой не е универсално завършен за всеки backend, но архитектурно показва, че библиотеката не се ограничава само до query execution. При достатъчно богата специализация на \code{connector\_trait<DB>} entity описанието може да участва и в еволюцията на схемата. Това променя и самия смисъл на ORM слоя: той вече не е само интерфейс към данните, а потенциален източник на schema-aware tooling.

\subsection{Архитектурни инварианти}

На базата на разгледаните слоеве могат да се формулират няколко инварианта. Първо, приложният код никога не бива да зависи пряко от драйверен API. Второ, query логиката се формулира в типизиран IR, а не в backend-native низове. Трето, разширяването към нов backend се извършва чрез trait специализация и capability декларации, а не чрез промяна на самия query API. Четвърто, ако една операция е недопустима за даден backend, това трябва да стане видимо възможно най-рано --- за предпочитане при компилация.

Тези инварианти правят архитектурата едновременно инженерно подредена и академично защитима. Те позволяват всяко следващо разширение да бъде оценявано не само по това дали ``работи'', а и по това дали запазва вече формулираните слоеве и граници.

\subsection{Архитектурен извод}

Обобщено, архитектурата на библиотеката разделя проблема на три равнища: логически модел, типизиран query IR и backend materialization. Това разделение позволява едновременно строга типова дисциплина, формализирана разширяемост и контролирано поведение върху различни модели за съхранение. В следващата глава това архитектурно решение се разглежда по-конкретно през реализацията на отделните подсистеми, където става ясно как архитектурният план се превръща в изпълним C++ код.
